\documentclass[11pt,a4paper]{article}
% Shared preamble for RuPhO English problem statements (match T1 2026 style).
% Use: \input{latex/rupho_en_preamble.tex} after \documentclass.
\usepackage[margin=1in]{geometry}
\usepackage{amsmath,amssymb}
\usepackage{graphicx}
\usepackage{float}
\usepackage{enumitem}
\usepackage{microtype}
\usepackage{caption}
\usepackage{tabularx}
\usepackage{hyperref}

\captionsetup{font=small,labelfont=bf,skip=6pt}

\setlist[itemize]{leftmargin=1.2cm}
\setlist[enumerate]{leftmargin=1.2cm}

% One outer box per subproblem: [id | statement | points] (no inner rules)
\newcommand{\problembox}[3]{%
  \par\addvspace{0.42em}%
  \noindent
  \setlength{\fboxsep}{7.5pt}%
  \setlength{\fboxrule}{0.95pt}%
  \fbox{%
    \begin{minipage}{\dimexpr\linewidth-2\fboxsep-2\fboxrule\relax}
    \begin{tabularx}{\linewidth}{@{}>{\raggedright\arraybackslash\bfseries}p{2.1em} @{\hspace{0.9em}} >{\raggedright\arraybackslash}X @{\hspace{0.9em}} >{\raggedleft\arraybackslash\bfseries}p{2.4em}@{}}
    #1 & #2 & #3
    \end{tabularx}
    \end{minipage}%
  }%
  \par\addvspace{0.32em}%
}

\title{\textbf{Double star}}
\author{\normalsize X17 (2017) --- English translation}
\date{}
\begin{document}
\maketitle

A binary star is a system of two gravitationally bound stars revolving around their center of mass in circular orbits. Cosmonaut-researcher Gluck wants to find such orbits that, when circling, his spacecraft would be on a straight line connecting these stars, and so that fuel consumption would be minimal. Use the following notations: $L$ – distance between stars, $L_x=xL$ – distance from the spacecraft to the star $M_1$, star mass ratio $k=M_1/M_2$.

\problembox{A1}{Express the parameter $k=k(x)$.}{4.50}

\problembox{A2}{Estimate numerically at what distances from the Moon Gluck’s ship must be in order for it to move in the Earth-Moon system as described in the problem statement. Consider the ratio of the masses of the Moon and the Earth to be $k=1/81$, and the distance from the Earth to the Moon to be $L=3{.}8\cdot{10^5~\text{km}}$.}{3.00}

\problembox{A3}{Schematically depict the possible positions of the satellite in the Earth-Moon system.}{0.50}

\vspace{1em}
\noindent\textit{Source:} \url{https://pho.rs/p/302}

\end{document}